\documentclass[11pt]{article}
\usepackage[margin=1in]{geometry}
\usepackage{amsmath,amssymb,amsthm}
\usepackage{graphicx}
\usepackage[hidelinks]{hyperref}

\newtheorem*{theorem}{Literature Answer}

\title{Literature Answer to Problem 1.1 of arXiv:1910.02414}
\author{FA/Banach run packet}
\date{2026-06-14}

\begin{document}
\maketitle

\section*{Status}

This is a literature-already-answered packet. It records that Problem 1.1
from Albiac--Ansorena--Dilworth--Kutzarova \cite{AADK} was answered
affirmatively in the later paper of Casazza--Dilworth--Kutzarova--Motakis
\cite{CDKM}. It is not an original proof from this run.

\section*{Original Problem Evidence}

The original problem appears as Problem 1.1 on page 4 of \cite{AADK}.

\begin{center}
\includegraphics[width=\textwidth]{figures/open_problem_crop.png}
\end{center}

In words, the problem asks whether there is a Banach space in which all
subsymmetric basic sequences are equivalent to a single basis, while that basis
is not symmetric.

\section*{Supporting Answer Evidence}

The later paper \cite{CDKM} explicitly cites the question from
\cite{KMP,AADK} and says that it answers it positively.

\begin{center}
\includegraphics[width=\textwidth]{figures/supporting_answer_crop.png}
\end{center}

\section*{Idea of the Answer}

The answer uses the subsymmetrization \(Su(T^*)\) of Tsirelson's original
space. The canonical basis of \(Su(T^*)\) is built to be subsymmetric but not
symmetric. The main rigidity result in the supporting paper shows that every
subsymmetric basic sequence in \(Su(T^*)\) is equivalent to this canonical
basis. Thus the basis supplies the required unique equivalence class, and its
failure of symmetry supplies the required non-symmetric witness.

\begin{theorem}
Problem 1.1 of \cite{AADK} has an affirmative answer. The Banach space
\(Su(T^*)\) from \cite{CDKM} has, up to equivalence, a unique subsymmetric
basic sequence, and its canonical subsymmetric basis is not symmetric.
\end{theorem}

\begin{proof}
Problem 1.1 in \cite{AADK} asks for a Banach space whose subsymmetric basic
sequences are all equivalent to one basis, with that basis non-symmetric.

The later paper \cite{CDKM} identifies this question in its introduction and
states that it answers the question positively. The witness is the
subsymmetrization \(Su(T^*)\) of Tsirelson's original space. In Section 2 of
\cite{CDKM}, the authors state that the canonical basis of \(Su(T^*)\) is not
symmetric, prove that the norm of \(Su(T^*)\) is \(1\)-subsymmetric, and prove
that every subsymmetric sequence in \(Su(T^*)\) is equivalent to the canonical
basis. These statements give exactly the existence requested in Problem 1.1.
\end{proof}

\section*{Limitations}

This packet answers only Problem 1.1 from \cite{AADK}. The adjacent Problems
1.2 and 1.3 ask for exactly two, or exactly \(n>1\), subsymmetric basic
sequences with at least one non-symmetric class. The supporting paper is not
used here as an answer to those finite-number questions.

\section*{Verification Notes}

The source and supporting papers are distinct. The supporting paper explicitly
mentions the question from \cite{AADK}, so this belongs in
\texttt{literature\_already\_answered} rather than
\texttt{literature\_implied\_answers}. The local duplicate search covered
arXiv ids \texttt{1910.02414} and \texttt{2112.09602}, plus the phrases
\texttt{unique subsymmetric}, \texttt{non-symmetric}, \texttt{Su(T*)}, and
\texttt{Tsirelson}.

\begin{thebibliography}{9}

\bibitem{AADK}
F. Albiac, J. L. Ansorena, S. J. Dilworth, and D. Kutzarova,
\emph{A dichotomy for subsymmetric basic sequences with applications to
Garling spaces}, Trans. Amer. Math. Soc. 374 (2021), no. 3, 2079--2106;
arXiv:1910.02414.

\bibitem{CDKM}
P. G. Casazza, S. J. Dilworth, D. Kutzarova, and P. Motakis,
\emph{On uniqueness and plentitude of subsymmetric sequences},
arXiv:2112.09602.

\bibitem{KMP}
D. Kutzarova, B. S. Mityagin, and K. Papathanasiou,
\emph{On subsymmetric basic sequences}, Studia Math. 218 (2013), no. 1,
1--21.

\end{thebibliography}

\end{document}
